% !PRINTED 598
\weberpage{598}
%<<CONT
seinerseits $R$ als Theiler enthält. Die Zahlen von $\Omega'$ bleiben
durch die Substitutionen $\Phi'$ ungeändert, und $\Phi'$ ist die Gruppe des
Körpers $\Omega$ in Bezug auf $\Omega'$ (vergl. Bd. I, \S.\,155). Wir bezeichnen
den Grad von $\Phi'$, der ein Theiler von $n$ ist, mit $n'$ und setzen

\weq{1}{n = m'\,n'.}

Es ist dann $n'$ der Grad von $\Omega$ in Bezug auf $\Omega'$, und $m'$
der Grad von $\Omega'$ in Bezug auf $R$.

Es ist dabei zu beachten, dass zwar $\Omega$ ein Normalkörper
auch in Bezug auf $\Omega'$ ist, dass aber $\Omega'$ im Allgemeinen kein
Normalkörper in Bezug auf $R$ sein wird.

Die Primideale im Körper $\Omega'$ lassen sich nun vollständig aus
denen von $\Omega$ ableiten.

Wenn $\mathfrak{p}$ irgend ein Primideal in $\Omega$ ist, so giebt es \gesperrt{ein und
nur ein durch $\mathfrak{p}$ theilbares Ideal $\mathfrak{p}'$ in $\Omega'$}, und $\mathfrak{p}'$ kann als
Theiler von $\mathfrak{P}$ durch keine anderen als die mit $\mathfrak{p}$ conjugirten Ideale
theilbar sein. Durchläuft $\varphi'$ die Substitutionen von $\Phi'$, so ist $\mathfrak{p}'$,
da es durch $\varphi'$ ungeändert bleibt, auch durch $\mathfrak{p}\,|\,\varphi'$ theilbar.

Wenn also $\mathfrak{p}_1'$ durch $\mathfrak{p}\,|\,\varphi_1$ theilbar ist, so ist es auch durch
$\mathfrak{p}\,|\,\psi\,\varphi_1\,\varphi'$ theilbar, wenn $\psi$ ein beliebiges Element der Gruppe
$\boldsymbol{\Psi}$ (\S.\,162) bedeutet.

Wenn daher $\varphi$ irgend ein Element des Systemes

\[
\Phi_1 = \boldsymbol{\Psi}\,\varphi_1\,\Phi'
\]

ist, so ist $\mathfrak{p}_1'$ durch $\mathfrak{p}\,|\,\varphi$ theilbar. Es kommt demnach jetzt die
Gruppenzerlegung in Betracht, die wir im \S.\,28 dieses Bandes
auseinandergesetzt haben.

Wenn $\varphi_2$ ein nicht in $\Phi_1$ enthaltenes Element aus $\Phi$ ist, so
hat das System

\[
\Phi_2 = \boldsymbol{\Psi}\,\varphi_2\,\Phi'
\]

mit $\Phi_1$ gar kein Element gemein. Giebt es noch ein Element $\varphi_3$,
das weder in $\Phi_1$ noch in $\Phi_2$ vorkommt, so bildet man ebenso

\[
\Phi_3 = \boldsymbol{\Psi}\,\varphi_3\,\Phi',
\]

und fährt so fort, bis die ganze Gruppe $\Phi$ erschöpft ist. Man
erhält so die Zerlegung

\weq{2}{\Phi = \Phi_1 + \Phi_2 + \Phi_3 + \cdots + \Phi_{e'}.}

Den Grad eines dieser Systeme $\Phi_r$ können wir nach \S.\,28 bestimmen.
Er ist gleich dem Producte des Grades von $\Phi$ multiplicirt mit % [sic] source prints $\Phi$; the equivalent expression below requires $\boldsymbol\Psi$
dem Index des Durchschnittes $\Psi_r'$ von $\Phi'$ mit $\Psi_r = \varphi_r^{-1}\boldsymbol{\Psi}\varphi_r$ in
Bezug auf $\Phi'$, oder was dasselbe ist, gleich dem Producte der
Grade von $\boldsymbol{\Psi}$ und $\Phi'$, getheilt durch den
%CONT>>
